% !PRINTED 464
\weberpage{464}
%<<CONT
\weq{10}{
\begin{aligned}
&\beta\partial = Ax,\\
&\alpha c - \gamma a = Cx,\\
&\beta c - \delta a + \alpha\partial = Bx,\\
&\beta c - \delta a - \alpha\partial = y,\\
&\beta c - \delta a = \frac{Bx+y}{2}, \quad \alpha\partial = \frac{Bx-y}{2},\\
&4n = y^2 - Dx^2.
\end{aligned}
}

Nehmen wir $A$ und $n$ ohne gemeinsamen Teiler an, so mu{\ss} $\partial = 1$, $a = n$ sein, denn eine in $x$ und $\partial$ aufgehende Primzahl m\"u{\ss}te auch in $y$ und folglich in $\alpha c - \gamma a$ und $\beta c - \delta a$ und, wegen $\alpha\delta - \beta\gamma = 1$, auch in $a$ und $c$ aufgehen, w\"ahrend doch $a$, $c$, $\partial$ keinen gemeinsamen Teiler haben sollen. Aus (10) wird $\alpha$, $\beta$, $\gamma$, $\delta$, $c$ ebenso bestimmt, wie im vorigen Paragraphen. Es ergeben sich zun\"achst die beiden Kongruenzen
\[
cAx \equiv \frac{Bx+y}{2}, \quad c\,\frac{Bx-y}{2} \equiv Cx \ (\mathrm{mod}\ n),
\]
von denen die eine aus der anderen folgt, wenn $x$ teilerfremd zu $n$ ist, und die Bedingung
\[
c \equiv 0 \ (\mathrm{mod}\ 2)
\]
ist damit vertr\"aglich. Weiter folgt:
\weq{11}{
\begin{aligned}
&\alpha = \frac{Bx-y}{2}, \quad n\gamma = -Cx + c\,\frac{Bx-y}{2},\\
&\beta = Ax, \qquad\qquad n\delta = -\frac{Bx+y}{2} + Acx.
\end{aligned}
}

Wenn nun $D$ und mithin $B$ ungerade ist, so m\"ussen, da wegen (9) $\alpha$ und $\delta$ beide gerade oder beide ungerade sind, $x$ und $y$ gerade sein; also sind $\beta$ und $\gamma$ gerade, und $\alpha$ und $\delta$ m\"ussen ungerade sein. Folglich mu{\ss} von den beiden Zahlen $\frac{1}{2}x$, $\frac{1}{2}y$ eine gerade, eine ungerade sein. Das ist aber unabh\"angig davon, wie sich $A$ und $C$ gegen den Modul 2 verhalten.

Wenn also von den drei Gr\"o{\ss}en (2) die eine der Gleichung (7) gen\"ugt, so tun es auch die beiden anderen, und wir k\"onnen diese drei nicht voneinander trennen.

Wir m\"ussen daher $B$ gerade annehmen und setzen zur Vereinfachung:
\weq{12}{D = -4m,}
worin $m$ eine positive ganze Zahl ist.

Wir unterscheiden:
%CONT>>
